\chapter{Sistemas (y se\~{n}ales)}

\section{Introducci\'{o}n}

\subsection{Definiciones}

\begin{enumerate}
\item Las se\~{n}ales son magnitudes naturales variables portadoras
  de informaci\'{o}n entre sistemas, representadas por funciones de
  variables independientes.

\item Los sistemas son conjuntos de partes que interact\'{u}an entre
  si, con un prop\'{o}sito com\'{u}n.

\item Los sistemas reciben se\~{n}ales de entrada para entregar
  se\~{n}ales de salida.

\item Sistema de tiempo continuo es aquel que recibe, procesa y
  entrega solamente se\~{n}ales de tiempo continuo.

\item Sistema de tiempo discreto es aquel que recibe, procesa y entrega
  solamente se\~{n}ales de tiempo discreto.

\item Sistema h\'{\i}brido es aquel que recibe, procesa y entrega se\~{n}ales
  de tiempo continuo y discreto.

\item Los sistemas pueden representarse mediante relaciones entre las
  se\~{n}ales de entrada y salida.

\item En el caso de sistemas de tiempo continuo, simbolizados
    \begin{equation}
      x(t) 
      \begin{array}{l}
        S               \\ 
        \rightarrow
      \end{array}
      y(t),
    \end{equation}
    podemos emplear los siguientes operadores matem\'{a}ticos para describirlos
    \begin{enumerate}
    \item Ecuaciones algebraicas
      \begin{equation}
        F(t, x(t), y(t)) = 0,
      \end{equation}
      donde $x(t)$ es la entrada e $y(t)$ es la salida del sistema. Por ejemplo,
      \begin{equation}
         y(t) - x(t) \, cos(2 \, t) = 0.
      \end{equation}
                                
    \item Ecuaciones diferenciales donde $x(t)$ es la entrada e $y(t)$ es la salida del sistema. Por ejemplo,
      \begin{equation}
        \frac{d y(t)}{dt} + 2 \, y(t) - x(t) = 0.
      \end{equation}
    \end{enumerate}
        
  \item En el caso de sistemas de tiempo discreto, simbolizados
    \begin{equation}
      x\left[ n\right] 
      \begin{array}{l}
        S               \\ 
        \rightarrow
      \end{array}
      y\left[ n\right] ,
    \end{equation}
    podemos emplear los siguientes operadores matem\'{a}ticos para describirlos
    \begin{enumerate}
    \item Ecuaciones algebraicas
      donde $x[n]$ es la entrada e $y[n]$ es la salida del sistema. Por ejemplo,
      \begin{equation}
         y[n] - x[n] \, 2^{-n} = 0.
      \end{equation}
                                
    \item Ecuaciones de diferencias finitas donde $x[n]$ es la entrada e $y[n]$ es la salida del sistema. Por ejemplo
      \begin{equation}
          y[n] - x[n] + x[n-1] - \frac{x[n-1]}{2} - y[n-1] - 2 \, y[n-2] + 4 \, y[n-3] = 0.
      \end{equation}
    \end{enumerate}

\item Las descripciones de sistemas, ya sean ecuaciones algebraicas,
  ecuaciones diferenciales, o de diferencias son modelos
  matem\'{a}ticos del sistema, y los podemos usar para an\'{a}lisis y
  dise\~{n}o.

\end{enumerate}

\subsection{Ejemplos}

\subsubsection{Divisor resistivo de tensi\'{o}n}
Supongamos que tenemos un divisor resistivo de tensi\'{o}n donde la
relaci\'{o}n entre la tensi\'{o}n de entrada $v_{e} \left( t \right)$
y salida $v_{s} \left( t \right)$ es la ecuaci\'{o}n
\begin{equation}
  V_{s}\left( t \right) = \frac{R_{2}}{R_{1} + R_{2}} \, V_{e}\left( t \right).      \label{divires}
\end{equation}
Esta es una ecuaci\'{o}n algebraica de tiempo continuo.  Ver Figura
(\ref{circuitoeps})

\subsubsection{Circuito RC Serie}
Supongamos que tenemos un circuito RC serie descripto por la
ecuaci\'{o}n diferencial
\begin{equation}
  \frac{dv_{c}\left( t\right) }{dt}+\frac{1}{R \, C} \, v_{c}\left( t\right) =
  \frac{1}{R \, C} \, v_{s}\left( t\right) . \label{rcserie}
\end{equation}
donde
\begin{enumerate}
\item $v_{c}(t)$ tensi\'{o}n en el capacitor,
\item $v_{s}\left( t\right) $ tensi\'{o}n en la fuente,
\item $R$ valor de la resistencia,
\item $C$ valor del capacitor.
\end{enumerate}
Ver Figura (\ref{circuitoeps}).

\subsubsection{Un negocio}
Supongamos que el due\~{n}o de un negocio repone todos los d\'{\i}as
la mercader\'{\i}a.  La ecuaci\'{o}n que describe la relaci\'{o}n
entre la cantidad de mercader\'{\i}a que entra en el dep\'{o}sito y el
dinero necesario para comprarla es
\begin{equation}
  t[n] = p \, (m - x[n]).
\end{equation}
donde
\begin{enumerate}
\item $x[n]$ es la cantidad de mercader\'{\i}a en el dep\'{o}sito,
\item $m$ la cantidad m\'{a}xima de mercader\'{\i}a que puede tener
  guardada,
\item $p$ el precio unitario,
\item $t[n]$ el pago por la compra del d\'{\i}a $n$.
\end{enumerate}
Esta es una ecuaci\'{o}n algebraica de tiempo discreto.

\subsubsection{Cuenta bancaria}
Supongamos que tenemos una cuenta bancaria en la que obtenemos cierto
inter\'{e}s mensual por mantener cierto dep\'{o}sito de dinero.  La
ecuaci\'{o}n de diferencias que nos dice cuanto dinero tendremos el
pr\'{o}ximo mes es
\begin{equation}
  y\left[n+1\right] = (1+c/100) \, y[n]+x[n],
\end{equation}
donde
\begin{enumerate}
\item $y[n+1]$ dinero en la cuenta el pr\'{o}ximo mes,
\item $y[n]$ dinero en la cuenta este mes,
\item $x[n]$ dinero depositado o extra\'{\i}do este mes,
\item $c$ inter\'{e}s.
\end{enumerate}
Esta es una ecuaci\'{o}n de diferencias y por lo tanto, de tiempo
discreto.

\subsubsection{C\'{a}lculo de ra\'{\i}ces de polin\'{o}mios}
Supongamos que tenemos que calcular \textbf{una} ra\'{\i}z del
polin\'{o}mio
\begin{equation}
  Q(x) = 0.
\end{equation}
Para obtener \textbf{una} soluci\'{o}n podemos escribir la ecuaci\'{o}n de
diferencias
\begin{equation}
  x_{k+1} = x_{k} - \left(\frac{Q(x)}{\frac{d Q(x)}{dx}}\right)_{x=x_{k}}.
\end{equation}
Si damos a $x_{0}$ un valor suficientemente cercano a \textbf{la}
ra\'{\i}z buscada del polin\'{o}mio $Q(x)$, obtendremos una secuencia
que converge r\'{a}pidamente. Por ejemplo, si consideramos el polinomio
\begin{equation}
Q(x) = x^{2} - 2 = 0, 
\end{equation}
calcularemos la raíz cuadrada de $2$. Por lo tanto, iteremos la ecuaci\'{o}n de diferencias
\begin{equation}
  x_{k+1} = \frac{2 + x^{2}_{k}}{2 x_{k}}, \qquad x_{0}   = 2,
\end{equation}
para obtener los primeros cinco valores de la secuencia
\begin{equation}
  \begin{array}{llllll}
    2 & 1.5 & 1.41667 & 1.41422 & 1.41421 & ...
  \end{array}
\end{equation}
que como vemos se acerca r\'{a}pidamente al resultado correcto.

\section{Modelos}

Necesitamos describir los sistemas con los que trabajamos para poder
analizarlos, verificar sus propiedades, o para dise\~{n}ar nuevos
sistemas. Como ya vimos, podemos usar ecuaciones diferenciales y
ecuaciones de diferencias.

\subsection{Ecuaciones diferenciales}

Consideremos una ecuaci\'{o}n diferencial
\begin{equation}
  \frac{d^{N} y(t)}{dt^{N}} + 
  \sum_{i=0}^{N-1} a_{i} \, \frac{d^{i} y(t)}{dt^{i}} = 
  \sum_{j=0}^{M} b_{i} \, \frac{d^{j} x(t)}{dt^{j}},
\end{equation}
donde $x(t)$ es la entrada, $y(t)$ es la salida, y los coeficientes
$a_{i}$ y $b_{j}$ son constantes reales, $N \geq 1$, $M \geq
0$, $N > M$. Por ejemplo,
\begin{equation}
\frac{dy^{3}(t)}{dt} + 3 \, \frac{dy^{2}(t)}{dt^{2}} 
+ 2 \, \frac{dy(t)}{dt}
+ y(t) 
= x(t).
\end{equation}
Las ecuaciones diferenciales pueden escribirse como un sistema de ecuaciones
diferenciales de primer orden denominadas ecuaciones de estado empleando 
variables de estado. 
Una posible definici\'{o}n de las variables de estado es la siguiente
\begin{eqnarray}
  v_{1}   &=& y(t),                                               \\
  v_{k+1} &=& \frac{dv_{k}}{dt} = \frac{d^{k}y}{dt^{k}},  \, 0 < k < N .
\end{eqnarray}
Con esta definici\'{o}n podemos escribir las ecuaciones de estado
\begin{eqnarray}
  \frac{dv_{k}}{dt} &=& v_{k+1},  \, 1 \leq k < N,  \\
  \frac{dv_{N}}{dt} &=& G(t, x(t), v_{1}, v_{2},...,v_{N}).
\end{eqnarray}
Por ejemplo, para
\begin{equation}
\frac{dy^{3}(t)}{dt} + 3 \, \frac{dy^{2}(t)}{dt^{2}} 
+ 2 \, \frac{dy(t)}{dt}
+ y(t) 
= x(t),
\end{equation}
definimos
\begin{equation}
v_{1} = y(t),
\end{equation}
\begin{equation}
\frac{dv_{1}}{dt} = v_{2} = \frac{dy(t)}{dt},
\end{equation}
\begin{equation}
\frac{dv_{2}}{dt} = v_{3} = \frac{dy^{2}(t)}{dt^{2}},
\end{equation}
\begin{equation}
\frac{dv_{3}}{dt} = \frac{dy^{3}(t)}{dt^{3}} = -v_{1} - 2 \, v_{2} - 3 \, v_{3} + x(t).
\end{equation}
En forma matricial
\begin{equation}
\left[ \begin{array}{c}
\frac{dv_{1}}{dt} \\
\frac{dv_{2}}{dt} \\
\frac{dv_{3}}{dt}
\end{array} \right] 
= \left[ \begin{array}{ccc}
0 & 1 & 0 \\
0 & 0 & 1 \\
-1 & -2 & -3
\end{array} \right] \cdot
\left[ \begin{array}{c}
v_{1}\\
v_{2}\\
v_{3}
\end{array} \right]
+
\left[ \begin{array}{c}
0 \\
0 \\
1
\end{array} \right] \, x(t).
\end{equation}
Las principales razones para emplear ecuaciones de estado son las
siguientes:
\begin{enumerate}
        
\item La teor\'{\i}a de ecuaciones diferenciales en esta notaci\'{o}n
  permite trabajar con sistemas de varias entradas y varias salidas.
        
\item Es muy sencillo programar computadoras para trabajar con esta
  notaci\'{o}n.
        
\item Esta notaci\'{o}n permite realizar simulaciones por computadora
  de grandes sistemas con muchas entradas y muchas salidas.
\end{enumerate}

\subsection{Ecuaciones de diferencias}

Consideremos una ecuaci\'{o}n de diferencias
\begin{equation}
  F(n, x[n], y[n-1],..., y[n-M+1]) = y[n]. 
\end{equation}
donde $x[n]$ es la entrada, $y[n]$ es la salida, y $M > 1$.
Por ejemplo
\begin{equation}
y[n+1] + \frac{1}{2} \, y[n] + \frac{1}{4} \, y[n-1] + \frac{1}{8} \, y[n-2] - x[n] = 0.
\end{equation}
Las ecuaciones de diferencias pueden
escribirse como un sistema de ecuaciones de diferencias de primer orden
denominadas ecuaciones de estado. Una posible definici\'{o}n de las
variables de estado es
\begin{eqnarray}
  v_{1}[n]   &=& y[n],                                               \\
  v_{k+1}[n] &=& v_{k}[n-1] = y[n-k],  0 < k < M .
\end{eqnarray}
Con esta definici\'{o}n podemos escribir las ecuaciones de estado
\begin{eqnarray}
  v_{k}[n+1] &=& v_{k+1}[n],  1 \leq k < M,  \\
  v_{M}[n+1] &=& G(n, x[n], v_{1}[n], v_{2}[n],...,v_{M}[n]).
\end{eqnarray}
Por ejemplo, considerando la ecuación de diferencias
\begin{equation}
y[n+1] + \frac{1}{2} \, y[n] + \frac{1}{4} \, y[n-1] + \frac{1}{8} \, y[n-2] - x[n] = 0,
\end{equation}
podemos definir las variables de estado
\begin{equation}
v_{1}[n] = y[n-2],
\end{equation}
\begin{equation}
v_{2}[n] = y[n-1] = v_{1}[n+1]
\end{equation}
\begin{equation}
v_{3}[n] = y[n] = v_{2}[n+1],
\end{equation}
\begin{equation}
v_{3}[n+1] = y[n+1] = - \frac{1}{2} \, v_{3}[n] - \frac{1}{4} \, v_{2}[n] - \frac{1}{8} \, v_{1}[n] + x[n].
\end{equation}
En forma matricial
\begin{equation}
\left[ \begin{array}{c}
v_{1}[n+1] \\
v_{2}[n+1] \\
v_{3}[n+1]
\end{array} \right] 
= \left[ \begin{array}{ccc}
0 & 1 & 0 \\
0 & 0 & 1 \\
-\frac{1}{8} & -\frac{1}{4} & -\frac{1}{2}
\end{array} \right] \cdot
\left[ \begin{array}{c}
v_{1}[n]\\
v_{2}[n]\\
v_{3}[n]
\end{array} \right]
+
\left[ \begin{array}{c}
0 \\
0 \\
1
\end{array} \right] \, x[n].
\end{equation}

\subsection{De ecuación diferencial a ecuación de diferencias}

\subsubsection{Aproximando la derivada, hacia adelante}

Sabemos que por definición de derivada
\begin{equation}
 \frac{dy}{dt} = \begin{array}{c}
                  Lim \\
                  \Delta \rightarrow 0
                 \end{array}
                 \frac{y(t+\Delta) - y(t)}{\Delta}.
\end{equation}
Entonces, si cambiamos el límite por un $\Delta$ pequeño, podemos escribir
\begin{equation}
 \frac{dy}{dt} \approx \frac{y[n+1] - y[n]}{\Delta}.
\end{equation}
Entonces, si tenemos la ecuación diferencial
\begin{equation}
 \frac{dy(t)}{dt} + 2 \, y(t) = x(t),
\end{equation}
resulta la ecuación de diferencias
\begin{equation}
 \frac{y[n+1] - y[n]}{\Delta} + 2 \, y[n] = x[n],
\end{equation}
o
\begin{equation}
(y[n+1] - y[n]) + 2 \,  \Delta \, y[n] = \Delta \, x[n],
\end{equation}
y finalmente
\begin{equation}
y[n+1] + (2 \, \times \Delta - 1) \, y[n] = \Delta \, x[n].
\end{equation}

\subsubsection{Aproximando la derivada, hacia atrás}

Podemos cambiar levemente la definición de derivada
\begin{equation}
 \frac{dy}{dt} = \begin{array}{c}
                  Lim \\
                  \Delta \rightarrow 0
                 \end{array}
                 \frac{y(t) - y(t-\Delta)}{\Delta}.
\end{equation}
Entonces, si cambiamos el límite por un $\Delta$ pequeño, podemos escribir
\begin{equation}
 \frac{dy}{dt} \approx \frac{y[n] - y[n-1]}{\Delta}.
\end{equation}
Entonces, si tenemos la ecuación diferencial
\begin{equation}
 \frac{dy(t)}{dt} + 2 \, y(t) = x(t),
\end{equation}
resulta la ecuación de diferencias
\begin{equation}
 \frac{y[n] - y[n-1]}{\Delta} + 2 \, y[n] = x[n],
\end{equation}
o
\begin{equation}
(y[n] - y[n-1]) + 2 \,  \Delta \, y[n] = \Delta \, x[n],
\end{equation}
y finalmente
\begin{equation}
(1 + 2 \, \times \Delta) \, y[n] - y[n-1] = \Delta \, x[n].
\end{equation}

\subsubsection{Otra aproximación de derivada}
Otra forma de definir derivada es
\begin{equation}
 \frac{dy}{dt} = \begin{array}{c}
                  Lim \\
                  \Delta \rightarrow 0
                 \end{array}
                 \frac{y(t+\Delta) - y(t-\Delta)}{2 \, \Delta}.
\end{equation}
Entonces podemos aproximar, para un $\Delta$ pequeño
\begin{equation}
 \frac{dy}{dt} \approx
                 \frac{y(t+\Delta) - y(t-\Delta)}{2 \, \Delta}.
\end{equation}
Entonces, si tenemos la ecuación diferencial
\begin{equation}
 \frac{dy(t)}{dt} + 2 \, y(t) = x(t),
\end{equation}
resulta la ecuación de diferencias
\begin{equation}
 \frac{y[n+1] - y[n-1]}{2 \, \Delta} + 2 \, y[n] = x[n],
\end{equation}
o
\begin{equation}
(y[n+1] - y[n-1]) + 4 \, \Delta \, y[n] = 2 \, \Delta \, x[n],
\end{equation}
y finalmente
\begin{equation}
y[n+1] + 4 \, \Delta \, y[n] - y[n-1] = 2 \, \Delta \, x[n].
\end{equation}


\section{Propiedades de sistemas}

Vamos a establecer un vocabulario b\'{a}sico para estudiar distintas
caracter\'{\i}sticas de sistemas.

\subsection{Sistemas sin memoria o est\'{a}ticos}

Los sistemas sin memoria o est\'{a}ticos son aquellos en los que la
salida actual depende \'{u}nicamente de la entrada actual.

\begin{equation}
  y\left( t\right) = 2 \, x\left(t\right) ,
\end{equation}
\begin{equation}
  y[n]=(x[n])^{3} + 4 \, x[n].
\end{equation}

\subsection{Sistemas con memoria o din\'{a}micos}

Los sistemas con memoria o din\'{a}micos son aquellos en los que la
salida actual depende de entradas de un momento distinto al actual.
Por ejemplo,
\begin{equation}
 \frac{dy(t)}{dt} + a \, y(t) = x(t),
\end{equation}
\begin{equation}
 \int_{t=0}^{T} x(t) \, dt = y(t),
\end{equation}
\begin{equation}
  y(t)=x(t)-x(t-\tau ),
\end{equation}
\begin{equation}
  y\left[ n\right] =\sum_{k=0}^{2}x[n-k]=x[n]+x[n-1]+x[n-2].
\end{equation}

\subsection{Sistemas inversos}

Un sistema $S$ tiene un sistema inverso $S^{-1}$ si su composici\'{o}n
es tal que
\begin{equation}
  x\left( t\right) =S^{-1}\left( S\left( x\left( t\right) \right) \right) ,
\end{equation}
\begin{equation}
  x\left[ n\right] =S^{-1}\left( S\left( x\left[ n\right] \right) \right) ,
\end{equation}

Tengamos en cuenta que un sistema puede no tener sistema inverso. Por
ejemplo
\begin{equation}
  y(t)=\left( x(t)\right) ^{2}.
\end{equation}

\subsection{Sistemas causales}

Un sistema es causal o no anticipativo si su salida no depende de
entradas futuras. Por ejemplo
\begin{equation}
  y\left( t\right) =ax\left( t\right),
\end{equation}
\begin{equation}
  y\left[ n\right] =x\left[ n-1\right] +x\left[ n-2\right] .
\end{equation}
Todos los sistemas sin memoria son causales.

\subsection{Sistemas no causales}

Un sistema es no causal o anticipativo si su salida depende de
entradas futuras. Por ejemplo, consideremos la ecuaci\'{o}n
\begin{equation}
  y\left[ n\right] =\frac{1}{5} \, \sum_{k=-2}^{2}x\left[ n-k\right] .
\end{equation}
Calculemos la salidad para $y[2]$. Expandiendo la sumatoria obtenemos
\begin{equation}
  y[2] = \frac{1}{5} \, \left( x[2-(-2)] + x[2-(-1)] + x[2-(0)] + x[2-(1)] + x[2-(2)] \right).
\end{equation}
Simplificando, resulta
\begin{equation}
  y[2] = \frac{1}{5} \left( x[4] + x[3] + x[2] + x[1] + x[0] \right).
\end{equation}
Como podemos ver, para calcular $y[2]$ necesitamos conocer las
entradas futuras $x[4]$ y $x[3]$. Por lo tanto, el sistema es no causal.

\subsection{Sistemas din\'{a}micos aut\'{o}nomos}

Un sistema din\'{a}mico es aut\'{o}nomo si
\begin{itemize}
\item Su ecuaci\'{o}n diferencial es del tipo
  \begin{equation}
    F(y(t), \frac{dy}{dt},..., \frac{d^{M}y}{dt^{M}}) = 0,                  
  \end{equation}
  donde $y(t)$ es la salida del sistema.

\item Su ecuaci\'{o}n de diferencias es del tipo
  \begin{equation}
    F(y[n], ..., y[n-M]) = 0, 
  \end{equation}
  donde $y[n]$ es la salida del sistema.
\end{itemize}

\subsection{Sistemas din\'{a}micos no aut\'{o}nomos}

Un sistema din\'{a}mico es no aut\'{o}nomo si
\begin{itemize}
\item Su ecuaci\'{o}n diferencial es del tipo
  \begin{equation}
    F(t, x(t), \frac{dx}{dt},..., \frac{d^{N}x}{dt^{N}}, y(t), \frac{dy}{dt},..., \frac{d^{M}y}{dt^{M}}) = 0,
  \end{equation}
  donde $x(t)$ es la entrada e $y(t)$ es la salida del sistema.

\item Su ecuaci\'{o}n de diferencias es del tipo
  \begin{equation}
    F(n, x[n], ..., x[n-N], y[n], ..., y[n-M]) = 0, 
  \end{equation}
  donde $x[n]$ es la entrada e $y[n]$ es la salida del sistema.
\end{itemize}

\subsection{Linealidad homogénea}

Una funci\'{o}n es lineal homogénea s\'{\i} y solo s\'{\i}
\begin{equation}
  f\left(a \, x_{1} + b \, x_{2}\right) = a \, f\left(x_{1}\right) + b \, f\left(x_{2}\right) .
\end{equation}
Por ejemplo,
\begin{equation}
  y = f(x) = m \, x,
\end{equation}
es una funci\'{o}n lineal, porque se cumple
\begin{equation}
  \begin{array}{ll}
    y_{1} & = m \, x_{1}, \\ 
    y_{2} & = m \, x_{2}, \\ 
    a \, y_{1}+ b \, y_{2} & = a \, m \, x_{1} + b \, m \, x_{2} = 
       m \, \left( a \, x_{1}+ b \, x_{2}\right) .
  \end{array}
\end{equation}
En cambio,
\begin{equation}
  y = f(x) = m \, x + 1,
\end{equation}
no es una funci\'{o}n lineal como podemos verificar
\begin{equation}
  \begin{array}{ll}
    y_{1} & = m \, x_{1}+1, \\ 
    y_{2} & = m \, x_{2}+1, \\ 
    a \, y_{1}+ b \, y_{2} & =
      m \, \left( a\, x_{1}+ b \, x_{2}\right) +\left( a+b\right) \neq
    m \, \left( a \, x_{1}+b \, x_{2}\right) +1.
  \end{array}
\end{equation}

\subsection{Sistemas lineales}

Se dice que un sistema $S$ es lineal si su respuesta para una entrada
compuesta $a \, x_{1}+b \, x_{2}$ cumple con la siguiente propiedad,
\begin{equation}
  S\left( a \, x_{1}+ b\, x_{2}\right) =a \, S\left( x_{1}\right) +b \, S\left( x_{2}\right),
\end{equation}
conocida como principio de superposici\'{o}n.

\subsection{Sistemas no lineales}

Se dice que un sistema $f$ es no lineal si su respuesta para una
entrada compuesta $a \, x_{1}+b \, x_{2}$ cumple con la siguiente propiedad,
\begin{equation}
  f\left( a \, x_{1}+b \, x_{2}\right) \neq a \, f\left( x_{1}\right) +b \, f\left( x_{2}\right) ,
\end{equation}
y por lo tanto, \textbf{NO} cumple con el principio de superposici\'{o}n.
Por ejemplo, esta ecuación diferencial no es lineal
\begin{equation}
 \frac{dy}{dt} = y(t) \, (M - y(t)) + x(t).
\end{equation}

\subsection{Puntos de equilibrio}

\begin{enumerate}

\item Un sistema din\'{a}mico de tiempo continuo con ecuación de estado
  \begin{equation}
    \mathbf{\dot x} = f(\mathbf{x}), 
  \end{equation}
tiene un punto de equilibrio $\mathbf{x_{e}}$ si se satisface la siguiente
  ecuaci\'{o}n
  \begin{equation}
    \mathbf{0} = f(\mathbf{x_{e}}).
  \end{equation}
        
\item Un sistema din\'{a}mico de tiempo discreto con ecuaci\'{o}n de estados
  \begin{equation}
    \mathbf{x_{k+1}} = f(\mathbf{x_{k}}), 
  \end{equation}
  tiene un punto de
  equilibrio $\mathbf{x_{e}}$ si se satisface la siguiente
  ecuaci\'{o}n
  \begin{equation}
    \mathbf{x_{e}} = f(\mathbf{x_{e}}).
  \end{equation}

\end{enumerate}
                
\subsection{Sistemas din\'{a}micos estables}

\begin{enumerate}
\item \textbf{Sistemas no aut\'{o}nomos} (\textbf{B}ounded
  \textbf{I}nput \textbf{B}ounded \textbf{O}utput) La salida permanece
  acotada para entradas acotadas.

\item \textbf{Sistemas aut\'{o}nomos} El sistema estable evoluciona
  por si s\'{o}lo a un estado de equilibrio.
\end{enumerate}

\subsection{Sistemas din\'{a}micos inestables}

Aquellos que no son estables. Apartados del equilibrio, se apartan
a\'{u}n m\'{a}s de \'{e}l.

\subsection{Sistemas invariantes en el tiempo}

Un sistema invariante en el tiempo es aqu\'{e}l cuyas caracter\'{\i}sticas 
y comportamiento est\'{a}n fijas en el tiempo. 

Este tipo de sistemas para una entrada determinada generan la misma
salida independientemente del momento.
\begin{enumerate}
\item La salida $y(t)$ para la entrada $x(t+t_{0})$ es independiente
  del valor $t_{0}$.
        
\item la salida $y[n]$ para la entrada $x[n + n_{0}]$ es independiente
  del valor $n_{0}$.
\end{enumerate}

\subsection{Sistemas variantes en el tiempo}


Si un sistema no es invariante en el tiempo, entonces se dice que es
variante.  En realidad, todos los sistemas f\'{\i}sicos \textbf{son}
variantes, debido al deterioro natural impuesto por la
termodin\'{a}mica. Lo que sucede es que con una escala de tiempo
adecuada, el deterioro puede ignorarse.

\subsection{Ejemplos}

\subsubsection{Clasificaci\'{o}n}

Supongamos un sistema definido por
\begin{equation}
  y(t)=g(t) \, x\left( t\right) .
\end{equation}

\textquestiondown Es lineal?  Si, como podemos ver en las siguientes
ecuaciones
\begin{eqnarray}
  y(t) &=& g(t) \, \left( a \, x_{1}(t)+b \, x_{2}(t)\right) =a \, g(t) \, x_{1}(t)+b \, g(t) \, x_{2}(t),    \\
  y(t) &=& a \, g(t) \, x_{1}(t)+b \, g(t) \, x_{2}(t)=a \, y_{1}(t)+b \, y_{2}(t).
\end{eqnarray}

\textquestiondown Es invariante en el tiempo?  No, como podemos ver en
las siguientes ecuaciones
\begin{eqnarray}
  y_{1}\left( t\right) &=&g\left( t\right) \, x_{1}\left( t\right) 
  \longleftrightarrow 
  y_{1}\left( t-\tau \right) =g\left( t-\tau \right) \, x_{1}\left( t-\tau \right) , \\
  y_{2}\left( t\right) &=&g\left( t\right) x_{1}\left( t-\tau \right) .
\end{eqnarray}

\textquestiondown Tiene memoria?  No, porque no necesitamos conocer
otro valor que el de la entrada actual para calcular la salida, ni
tampoco sus derivadas.

\subsubsection{Un circuito el\'{e}ctrico}

Supongamos un circuito el\'{e}ctrico con las siguientes magnitudes:
corriente $i(t),$ tensi\'{o}n $v(t),$ resistencia $R,$ capacidad $C$,
inductancia $L$.
\begin{equation}
  i\left( t\right) =
  C\frac{dv\left( t\right) }{dt}+
  \frac{v\left( t\right) }{R}+
  \frac{1}{L}\int_{0}^{t}v\left( \tau \right) d\tau .             \label{rlc}
\end{equation}
Ver Figura (\ref{circuitoeps}).

\textquestiondown Es lineal?  Si, porque las operaciones derivada e
integral son lineales
\begin{eqnarray}
  \frac{d(a x(t) + b y(t)}{dt} &=& a \frac{dx}{dt} + b \frac{dy}{dt},     \\
  \int a x(t) + b y(t) dt &=& a \int x(t) dt + b \int y(t) dt.
\end{eqnarray}

\textquestiondown Es invariante en el tiempo?  El sistema es
invariante en el tiempo. \textquestiondown Porqu\'{e}?

\textquestiondown Tiene memoria?  Si, por las operaciones derivada e
integral.

\subsubsection{Un sistema abstracto, tiempo continuo}

El sistema
\begin{equation}
  y(t)=S\left( x\left( t\right) \right) =2x\left( t\right) ,
\end{equation}
es un sistema sin memoria o est\'{a}tico de tiempo continuo que tiene
un sistema inverso
\begin{equation}
  S^{-1}\left( x\left( t\right) \right) =\frac{1}{2}x(t),
\end{equation}

\subsubsection{Otro sistema abstracto, tiempo discreto}

El sistema
\begin{equation}
  y\left[ n\right] =\sum_{k=0}^{\infty }x[n-k],
\end{equation}
es un sistema con memoria o din\'{a}mico de tiempo discreto que tiene
un sistema inverso
\begin{equation}
  x\left[ n\right] =y\left[ n\right] -y\left[ n-1\right] =
  \left(\sum_{k=0}^{\infty }x[n-k]\right) -\left( \sum_{k=0}^{\infty}x[n-1-k]\right) .
\end{equation}
\begin{equation}
  x\left[ n\right] =
  x[n]+\left( \sum_{k=1}^{\infty }x[n-k]\right) -\left(\sum_{k=0}^{\infty }x[n-1-k]\right) .
\end{equation}
\begin{equation}
  x\left[ n\right] =
  x[n]+\left( \sum_{p=0}^{\infty }x[n-1-p]\right) -\left(\sum_{k=0}^{\infty }x[n-1-k]\right) .
\end{equation}
\begin{equation}
  x\left[ n\right] =x[n].
\end{equation}

\begin{figure}
  \centering
  \includegraphics[totalheight=3.7in,trim=0 350 0
  0,clip]{./graphics/circuitos.eps}
  \caption{Distintos circuitos usados para plantear las ecua-ciones
    (\ref{divires}), (\ref{rcserie}) y (\ref{rlc})}
  \label{circuitoeps}
\end{figure}

\subsubsection{Sistema mec\'{a}nico}

Supongamos un sistema mec\'{a}nico con las siguientes magnitudes: masa
$M,$ constante de fricci\'{o}n $B,$ constante del resorte $K,$ fuerza
externa $f(t),$ velocidad de la masa $v(t)$.
\begin{equation}
  f(t)=M\frac{dv\left( t\right) }{dt}+Bv\left( t\right) +K\int_{-\infty}^{t}v\left( \tau \right) d\tau .
\end{equation}
\begin{enumerate}

\item \textquestiondown Es lineal?

\item \textquestiondown Es invariante en el tiempo?

\item \textquestiondown Tiene memoria?

\end{enumerate}

\subsubsection{Analog\'{\i}a entre diferentes sistemas}

Compare las ecuaciones diferenciales de los sistemas:
\begin{equation}
  i\left( t\right) =
  C\frac{dv\left( t\right) }{dt}+\frac{v\left( t\right) }{R}+
  \frac{1}{L}\int_{-\infty }^{t}v\left( \tau \right) d\tau ,
\end{equation}
\begin{equation}
  f(t)=
  M\frac{dv\left( t\right) }{dt}+Bv\left( t\right) +K\int_{-\infty}^{t}v\left( \tau \right) d\tau .
\end{equation}
Estas ecuaciones diferenciales corresponden a distintos sistemas (uno
el\'{e}ctrico, otro mec\'{a}nico), pero las dos tienen la misma forma.

\section{Diagrama de bloques}

\begin{enumerate}
\item Es posible hacer descripciones de sistemas usando
  \begin{enumerate}
  \item bloques que representan los distintos componentes del sistema.
  \item flechas que representan los flujos de se\~{n}al entre los
    componentes (que salida se conecta a cada entrada).
  \end{enumerate}
                  
\item Se pueden hacer diagramas de bloques del mismo sistema con
  \'{e}nfasis en distintos aspectos:
  \begin{enumerate}
  \item Caja negra, con \'{e}nfasis en los conjuntos de entradas y
    salidas (o interfases de entrada y salida). Vea la Fig.
    (\ref{cajanegra}).
        
  \item Caja blanca, con \'{e}nfasis en las relaciones entre los
    distintos componentes del sistema (o subsistemas). Vea la Fig.
    (\ref{cajablanca}).
        
  \end{enumerate}
  \begin{figure}
    \centering
    \includegraphics[height=4in,angle=-90]{./graphics/cajanegra.eps}
    \caption{Sistema visto como caja negra, con \'{e}nfasis en las
      interfaces de entrada y salida}
    \label{cajanegra}
  \end{figure}
  \begin{figure}
    \centering
    \includegraphics[height=4in,angle=-90]{./graphics/cajablanca.eps}
    \caption{Sistema visto como caja blanca, con \'{e}nfasis en los
      subsistemas}
    \label{cajablanca}
  \end{figure}

\item Existen varias formas de interconectar sistemas (Vea la Fig.
  \ref{conectioeps})
  \begin{enumerate}
  \item Conexión serie
        
  \item Conexión paralelo
                        
  \item Retroalimentación
        
  \end{enumerate}
  \begin{figure}
    \centering
    \includegraphics[width=4.5in]{./graphics/conectio.eps}
    \caption{Interconecci\'{o}n de sistemas. El bloque circular con
      dos signos $+$ en su interior es un sumador.}
    \label{conectioeps}
  \end{figure}

\end{enumerate}

\section{Ejemplos}

\subsection{Planteo}

Indique si los sistemas definidos por las siguientes relaciones entrada salida son lineales o no lineales
\begin{enumerate}
\item $y\left[ n+1\right] =\frac{1}{2}y\left[ n\right] +x\left[
    n\right]$,

\item $y\left[ n+1\right] =-\frac{1}{4}y\left[ n\right] +x\left[
    n\right]$,

\item $y\left[ n+1\right] =\left\{
    \begin{array}{c}
      \frac{1}{2}y\left[ n\right] +x\left[ n\right] , \\ 
      -\frac{1}{4}y\left[ n\right] +x\left[ n\right] ,
    \end{array}
    \begin{array}{c}
      y\left[ n\right] >\frac{1}{2}, \\ 
      y\left[ n\right] \leq \frac{1}{2}.
    \end{array}
  \right. $
\end{enumerate}

\subsection{Solución}

\begin{enumerate}
	\item $y\left[ n+1\right] =\frac{1}{2}y\left[ n\right] +x\left[
	n\right]$, Sistema lineal, ecuación de diferencias donde aparecen las señales multiplicadas por coeficientes constantes y sumadas.
	
	\item $y\left[ n+1\right] =-\frac{1}{4}y\left[ n\right] +x\left[
	n\right]$, Sistema lineal, ecuación de diferencias donde aparecen las señales multiplicadas por coeficientes constantes y sumadas.
	
	\item $y\left[ n+1\right] =\left\{
	\begin{array}{c}
		\frac{1}{2}y\left[ n\right] +x\left[ n\right] , \\ 
		-\frac{1}{4}y\left[ n\right] +x\left[ n\right] ,
	\end{array}
	\begin{array}{c}
		y\left[ n\right] >\frac{1}{2}, \\ 
		y\left[ n\right] \leq \frac{1}{2}.
	\end{array}
	\right. $ Sistema no lineal, o lineal a tramos. No cumple con el principio de superposición.
\end{enumerate}

\section{Para practicar}

Escriba un programa que calcule la ra\'{\i}z c\'{u}bica de $10$ mediante una ecuaci\'{o}n de diferencias. Haga la verificaci\'{o}n correspondiente.

Si planteamos
\begin{equation}
	p(x) = x^{3} - 10,
\end{equation}
podemos calcular la derivada
\begin{equation}
	\frac{dp(x)}{dx} = 3 \, x^{2}
\end{equation}
Entonces podemos escribir la ecuación de diferencias
\begin{equation}
	x[n+1] = x[n] - \frac{x[n]^{3} - 10}{3 \, x[n]^{2}},
\end{equation}
Supongamos $x[0] = 1$.


